\documentclass[aspectratio=169, 11pt, xcolor={table}]{beamer}
% Shared style for the algorithmics course slides (same template as the
% ENSB generative-AI course), with C code listings added.
% Dark blue removed: every slide (cover included) uses the light background
% of the template's content slides, with teal + orange as accents.

\usepackage{amsmath}
\usepackage{amssymb}
\usepackage{booktabs}
\usepackage{colortbl}
\usepackage{tabularx}
\usepackage{tikz}
\usepackage[skins]{tcolorbox}
\usepackage{fontawesome5}
\usepackage{listings}
\usepackage{polyglossia}

\usetikzlibrary{positioning, calc, arrows.meta, shapes.geometric}

%------------------------------------------------
% Languages and fonts (compile from slides/ with XeLaTeX or tectonic)
%------------------------------------------------
\setdefaultlanguage[numerals=maghrib]{arabic}
\setotherlanguage{english}
\usefonttheme{serif}

\newfontfamily\arabicfont[
    Script=Arabic, Scale=0.98, Path=../../fonts/, Extension=.ttf,
    UprightFont=*-Regular, BoldFont=*-Bold,
    ItalicFont=*-Italic, SlantedFont=*-Italic,
    BoldItalicFont=*-BoldItalic, BoldSlantedFont=*-BoldItalic]{Amiri}
\newfontfamily\arabicfontsf[
    Script=Arabic, Scale=0.98, Path=../../fonts/, Extension=.ttf,
    UprightFont=*-Regular, BoldFont=*-Bold,
    ItalicFont=*-Italic, SlantedFont=*-Italic,
    BoldItalicFont=*-BoldItalic, BoldSlantedFont=*-BoldItalic]{Amiri}
% Arabic fallback for \ttfamily (needed by the C listings)
\newfontfamily\arabicfonttt[Script=Arabic, Scale=0.98, Path=../../fonts/,
    Extension=.ttf, UprightFont=*-Regular, BoldFont=*-Bold,
    ItalicFont=*-Italic, SlantedFont=*-Italic,
    BoldItalicFont=*-BoldItalic, BoldSlantedFont=*-BoldItalic]{Amiri}
% Latin monospace font for C code (selected directly, so polyglossia
% does not swap it for the Arabic fallback above)
\newfontfamily\codefont[Scale=0.95, BoldFont=lmmonolt10-bold.otf,
    ItalicFont=lmmono10-italic.otf]{lmmono10-regular.otf}
% Heading font
\newfontfamily\khatt[Script=Arabic, Scale=1.05, Path=../../fonts/,
    BoldFont=ae_AlBattar.ttf]{ae_AlBattar.ttf}

%------------------------------------------------
% Colors (template palette without the navy tones)
%------------------------------------------------
\definecolor{brandTeal}{HTML}{5CC8E6}
\definecolor{brandOrange}{HTML}{FFA27F}
\definecolor{brandLight}{HTML}{F8F9FA}
\definecolor{brandDark}{HTML}{263238}
\definecolor{brandMuted}{HTML}{6C757D}
\colorlet{tealDeep}{brandTeal!60!black}
\definecolor{codeGrey}{HTML}{F1F3F5}

\setbeamercolor{background canvas}{bg=brandLight}
\setbeamercolor{normal text}{fg=brandDark}
\setbeamercolor{structure}{fg=brandTeal}
\setbeamertemplate{navigation symbols}{}

% beamer.cls sets \raggedright globally; under RTL that centers short lines
% or pushes them left, so go back to justified text (flush right).
\rightskip=0pt\relax
% Lists call \raggedright too; make it a no-op for RTL slides.
\AtBeginDocument{\renewcommand{\raggedright}{\rightskip=0pt\relax}}

%------------------------------------------------
% Frame title, footer, lists
%------------------------------------------------
\setbeamertemplate{frametitle}{%
    \vspace{0.35cm}
    \noindent{\khatt\bfseries\Large\color{brandDark}\insertframetitle}
    \par
    \vskip-2pt
    \noindent{\color{brandTeal}\rule{\textwidth}{1.5pt}}
    \vspace{0.05cm}
}

% Footer: page number only
\setbeamertemplate{footline}{%
    \hbox to \paperwidth{\hfill
        {\color{brandMuted}\scriptsize\LR{\insertframenumber\,/\,\inserttotalframenumber}}%
        \hspace{0.6cm}}%
    \vspace{0.25cm}}

\setbeamertemplate{itemize item}{\color{brandTeal}\small\faChevronLeft}
\setbeamertemplate{itemize subitem}{\color{brandMuted}\small\faAngleLeft}
\setbeamertemplate{enumerate item}{\color{tealDeep}\bfseries\insertenumlabel)}

% Tables: every cell framed by a visible border
\newcolumntype{R}{>{\raggedleft\arraybackslash}X}
\newcolumntype{C}{>{\centering\arraybackslash}X}
\newcommand{\ras}[1]{\cellcolor{brandTeal!20}\textbf{#1}}
% Colored text inside a table cell (\leavevmode keeps it on the row baseline)
\newcommand{\cell}[2]{\leavevmode\textcolor{#1}{#2}}
\renewcommand{\arraystretch}{1.35}
\setlength{\arrayrulewidth}{0.6pt}
\arrayrulecolor{brandMuted!70}

%------------------------------------------------
% Boxes (the two boxes of the template's 2nd slide)
%------------------------------------------------
\tcbset{
    brandbox/.style={
        enhanced, arc=4pt, outer arc=4pt, boxrule=1.5pt,
        coltitle=white, coltext=brandDark, fonttitle=\khatt\bfseries,
        toptitle=0.8mm, bottomtitle=0.8mm,
        top=1.2mm, bottom=1.2mm, left=3mm, right=3mm,
        before skip=4pt, after skip=4pt}
}
% Teal box: definitions, objectives, summaries
\newtcolorbox{tbox}[2][]{brandbox, colback=brandTeal!10,
    colframe=brandTeal, title={#2}, #1}
% Orange box: warnings
\newtcolorbox{obox}[2][]{brandbox, colback=brandOrange!10,
    colframe=brandOrange, title={#2}, #1}
% Card (template's 3rd slide, in teal)
\newtcolorbox{card}[2][]{enhanced, colback=brandTeal!6, colframe=brandTeal,
    arc=6pt, boxrule=1.2pt, halign=center, halign title=center,
    coltitle=white, coltext=brandDark, fonttitle=\khatt\bfseries,
    top=1.5mm, bottom=1.5mm, left=2mm, right=2mm, title={#2}, #1}
% Prompt text written to the tool
\newtcolorbox{amr}[1][brandTeal]{enhanced, arc=1mm, boxrule=0pt,
    colback=white, overlay={\fill[#1] (frame.north east) rectangle ([xshift=-2.5pt]frame.south east);},
    left=6pt, right=8pt, top=2pt, bottom=2pt,
    before skip=3pt, after skip=3pt}

%------------------------------------------------
% TikZ (Arabic right-to-left inside nodes)
%------------------------------------------------
\newcommand{\nodecolor}{}
\tikzset{
    every node/.append style={
        execute at begin node={\nodecolor\everypar{\beginR}\beginR},
        execute at end node=\endR},
    lawn/.code={\def\nodecolor{\color{#1}}},
    khatwa/.style={draw=#1, fill=#1!12, very thick, rounded corners=5pt,
        minimum height=0.95cm, text width=1.9cm, align=center,
        font=\khatt\bfseries\small},
    khatwa/.default=brandTeal,
    sahm/.style={-{Stealth[length=2.5mm]}, thick, color=brandMuted},
    rang/.style={draw=#1, fill=#1!12},
    raqm/.style={circle, fill=#1, lawn=white, inner sep=1.5pt,
        font=\small\bfseries},
}

% Horizontal chain of steps, right to left:
% \flow[width]{{a}, {b\\c}, {d}}
\newcommand{\flow}[2][1.9cm]{%
    \begin{tikzpicture}[node distance=5mm]
        \foreach \t [count=\i] in {#2} {
            \ifnum\i=1
                \node[khatwa, text width=#1] (f\i) {\t};
            \else
                \pgfmathtruncatemacro\p{\i-1}
                \node[khatwa, text width=#1, left=of f\p] (f\i) {\t};
                \draw[sahm] (f\p) -- (f\i);
            \fi
        }
    \end{tikzpicture}}

%------------------------------------------------
% C code (always left-to-right, inside a fragile frame)
%------------------------------------------------
\lstdefinestyle{cstyle}{
    language=C, basicstyle=\codefont\footnotesize, columns=fullflexible,
    keepspaces=true, showstringspaces=false, upquote=true,
    keywordstyle=\color{tealDeep}\bfseries,
    commentstyle=\color{red!75!black},
    stringstyle=\color{brandOrange!85!black},
    xleftmargin=16pt,
    numbers=left, numberstyle=\tiny\color{brandMuted}, numbersep=7pt,
    aboveskip=0pt, belowskip=0pt,
    morekeywords={bool, true, false}}
% Inline listings must sit in the english environment, not LTR, otherwise
% numbers such as 3.14 are reversed.
% Numbers in Arabic text: \textenglish{1.75} (not \LR).
% Compact code size for long programs (fixed line height, no gaps)
\newcommand{\codesmall}{\fontsize{8}{8}\selectfont}
\newcommand{\codetiny}{\fontsize{7.4}{7.6}\selectfont}
% Inline C: \code{printf}
\newcommand{\code}[1]{\mbox{\textenglish{\codefont\begingroup\color{tealDeep}#1\endgroup}}}
% Program output (console)
\lstdefinestyle{outstyle}{
    language={}, basicstyle=\codefont\footnotesize\color{brandDark},
    columns=fullflexible, keepspaces=true, showstringspaces=false,
    aboveskip=0pt, belowskip=0pt}

% Listings read from files (no fragile frame needed):
%   \ccode{d1/circle.c}   \out{d1/circle.out}
% One box per listing for the grey background (listings' own per-line
% background leaves thin white gaps between lines)
\newtcolorbox{codebox}{enhanced, colback=codeGrey, colframe=codeGrey,
    boxrule=0pt, arc=0pt, left=4pt, right=4pt, top=1pt, bottom=1pt,
    before skip=0pt, after skip=2pt,
    overlay={\fill[brandMuted!45] (frame.south west) rectangle ([xshift=2.5pt]frame.north west);}}
\newtcolorbox{outbox}{enhanced, colback=brandMuted!20, colframe=brandMuted!20,
    boxrule=0pt, arc=0pt, left=4pt, right=4pt, top=1.5pt, bottom=1.5pt,
    before skip=2pt, after skip=2pt}
\newcommand{\ccode}[2][]{\begin{codebox}\begin{english}\lstinputlisting[style=cstyle,#1]{../code/#2}\end{english}\end{codebox}}
\newcommand{\out}[2][]{\begin{outbox}\begin{english}\lstinputlisting[style=outstyle,#1]{../code/#2}\end{english}\end{outbox}}
% Audience line of the cover
\newcommand{\audience}{طلبة السنة الأولى}

%------------------------------------------------
% Cover (template's 1st slide on the light background)
%------------------------------------------------
\newcommand{\cover}[3]{%
\begin{frame}[plain]
    \begin{tikzpicture}[remember picture, overlay]
        \draw[brandTeal, line width=1.5pt, opacity=0.3]
            (current page.south east) ++(-1,1) circle (4cm);
        \draw[brandTeal, line width=0.8pt, opacity=0.25]
            (current page.north west) ++(2,2) circle (6cm);
        \draw[brandOrange, line width=0.5pt, opacity=0.25]
            (current page.north west) ++(2,2) circle (7cm);
        \node[anchor=north east] at ($(current page.north east)+(-1.0,-0.8)$) {
            \begin{tikzpicture}[scale=0.35]
                \fill[brandTeal] (0,0) -- (0.8,1.2) -- (1.6,0) -- cycle;
                \fill[brandOrange] (0.8,1.2) -- (1.6,2.4) -- (2.4,1.2) -- cycle;
            \end{tikzpicture}
        };
    \end{tikzpicture}

    \vspace{1.6cm}
    \begin{minipage}{0.85\textwidth}
        {\color{tealDeep}\khatt\bfseries #1} \\[0.3cm]
        {\raggedleft\color{brandDark}\Huge\khatt\bfseries #2\par} \vspace{0.3cm}
        {\color{brandMuted}\large #3}
    \end{minipage}

    \vspace{1.2cm}
    \begin{minipage}{0.85\textwidth}
        \color{brandMuted}\small
        \begin{tabular}{@{}ll}
            \textbf{الأستاذ:} & د. بن عبد الرزاق يوسف \\
            \textbf{الفئة:}   & \audience \\
        \end{tabular}
    \end{minipage}
\end{frame}}

\title{مدخل إلى الخوارزميات}

\begin{document}

\cover{درس الفصل الأول}
{مدخل إلى الخوارزميات}
{المتغيرات والثوابت والعمليات الأساسية بلغة \textenglish{C}}

\begin{frame}{أهداف الدرس}
    \begin{tbox}{\faBullseye\quad في نهاية هذا الدرس، تستطيع أن:}
        \begin{itemize}
            \item تعرّف الخوارزمية وتكتب هيكل برنامج بلغة \textenglish{C}.
            \item تصرّح بالمتغيرات والثوابت بلغة \textenglish{C}.
            \item تستعمل التعيين والقراءة والكتابة والعمليات الحسابية.
        \end{itemize}
    \end{tbox}
\end{frame}

\begin{frame}{ما الخوارزمية؟}
    \begin{tbox}{\faLightbulb\quad تعريف}
        \textbf{الخوارزمية}: مجموعة منتهية ومرتبة من التعليمات الواضحة،
        تحل مشكلة انطلاقًا من \textbf{معطيات} للوصول إلى \textbf{نتائج}.
    \end{tbox}
    يُنسب اسمها إلى العالم المسلم \textbf{محمد بن موسى الخوارزمي}.
\end{frame}

\begin{frame}{خصائص الخوارزمية}
    \vspace{0.6cm}
    \begin{columns}[T]
        \column{0.3\textwidth}
        \begin{card}{منتهية}
            تتوقف بعد عدد محدود من الخطوات.
        \end{card}
        \column{0.3\textwidth}
        \begin{card}{دقيقة}
            لكل تعليمة معنى واحد فقط.
        \end{card}
        \column{0.3\textwidth}
        \begin{card}{صحيحة}
            تعطي النتيجة الصحيحة دائمًا.
        \end{card}
    \end{columns}
\end{frame}

\begin{frame}{من المشكلة إلى البرنامج}
    \vspace{0.8cm}
    \begin{center}
        \flow[2.0cm]{{المشكلة}, {التحليل}, {الخوارزمية}, {البرنامج\\بلغة \textenglish{C}}, {التنفيذ}}
    \end{center}
    \vspace{0.4cm}
    نحلل المشكلة، ثم نكتب الخوارزمية، ثم نترجمها إلى برنامج يفهمه الحاسوب.
\end{frame}

\begin{frame}{هيكل برنامج بلغة \textenglish{C}}
    \begin{columns}[T]
        \column{0.45\textwidth}
        \ccode{d1/skeleton.c}
        \column{0.5\textwidth}
        \begin{itemize}
            \item \code{stdio.h}: مكتبة الإدخال والإخراج.
            \item \code{main}: نقطة بداية البرنامج.
            \item \code{return 0}: انتهى البرنامج بنجاح.
        \end{itemize}
    \end{columns}
\end{frame}

\begin{frame}{المتغير}
    \begin{tbox}{\faLightbulb\quad تعريف}
        \textbf{المتغير}: مكان في الذاكرة له \textbf{اسم} و\textbf{نوع} و\textbf{قيمة}
        يمكن أن تتغير أثناء التنفيذ.
    \end{tbox}
\end{frame}

\begin{frame}{المتغير في الذاكرة}
    \vspace{0.5cm}
    \begin{center}
    \begin{tikzpicture}[
        khana/.style={draw=brandTeal, fill=white, very thick, minimum width=2.2cm,
            minimum height=1cm, font=\Large\bfseries},
        lafta/.style={font=\small, lawn=brandMuted}]
        \node[khana] (a) at (0,0) {\textenglish{19}};
        \node[khana] (b) at (-3.2,0) {\textenglish{1.75}};
        \node[khana] (c) at (-6.4,0) {\textenglish{'A'}};
        \node[above=2pt of a] {\code{age}};
        \node[above=2pt of b] {\code{height}};
        \node[above=2pt of c] {\code{letter}};
        \node[lafta, below=2pt of a] {عدد صحيح};
        \node[lafta, below=2pt of b] {عدد حقيقي};
        \node[lafta, below=2pt of c] {حرف};
    \end{tikzpicture}
    \end{center}
    \vspace{0.3cm}
    المتغير \textbf{علبة} لها اسم. عند وضع قيمة جديدة فيها، \textbf{تُمحى القيمة القديمة}.
\end{frame}

\begin{frame}{أنواع المعطيات الأساسية}
    \begin{center}
    \begin{tabular}{|c|c|c|c|}
        \hline
        \ras{النوع} & \ras{\textenglish{C}} & \ras{الصيغة} & \ras{مثال} \\ \hline
        عدد صحيح  & \code{int}    & \code{\%d}  & \textenglish{-3, 25} \\ \hline
        عدد حقيقي & \code{float}  & \code{\%f}  & \textenglish{3.14} \\ \hline
        حرف       & \code{char}   & \code{\%c}  & \textenglish{'a'} \\ \hline
        منطقي     & \code{int}    & \code{\%d}  & \textenglish{0, 1} \\ \hline
    \end{tabular}
    \end{center}
    في لغة \textenglish{C}: القيمة \textenglish{0} تعني \textbf{خطأ}، وكل قيمة أخرى تعني \textbf{صحيح}.
\end{frame}

\begin{frame}{التصريح بالمتغيرات}
    \begin{center}
    \begin{minipage}{0.62\textwidth}
        \ccode{d1/decl.c}
    \end{minipage}
    \end{center}
    \vspace{0.3cm}
    في \textenglish{C}: \textbf{النوع} ثم \textbf{الاسم} ثم الفاصلة المنقوطة \code{;}.
\end{frame}

\begin{frame}{قواعد تسمية المتغيرات}
    \begin{columns}[T]
        \column{0.47\textwidth}
        \begin{tbox}{\faCheckCircle\quad صحيح}
            \code{age}\quad \code{note1}\quad \code{total\_price}
        \end{tbox}
        \column{0.47\textwidth}
        \begin{obox}{\faTimesCircle\quad خطأ}
            \code{1note}\quad \code{total price}\quad \code{int}
        \end{obox}
    \end{columns}
    \vspace{0.3cm}
    يبدأ الاسم بحرف لاتيني، دون فراغ، وليس كلمة محجوزة.\\
    \code{Age} و\code{age} متغيران مختلفان.
\end{frame}

\begin{frame}{الثوابت}
    \begin{center}
    \begin{minipage}{0.62\textwidth}
        \ccode{d1/const.c}
    \end{minipage}
    \end{center}
    \vspace{0.3cm}
    \textbf{الثابت}: قيمة لها اسم و\textbf{لا تتغير} أثناء التنفيذ.
\end{frame}

\begin{frame}{التعيين}
    \begin{center}
    \begin{minipage}{0.62\textwidth}
        \ccode{d1/assign.c}
    \end{minipage}
    \end{center}
    \vspace{0.3cm}
    نحسب الطرف الأيمن أولًا، ثم نضع النتيجة في المتغير.\\
    \code{x = x + 1} ليست معادلة، بل تعني: زِد \code{x} بواحد.
\end{frame}

\begin{frame}{اختصارات التعيين في \textenglish{C}}
    \begin{columns}[T]
        \column{0.5\textwidth}
        \ccode{d1/short.c}
        \column{0.45\textwidth}
        هذه الاختصارات كثيرة الاستعمال، خاصة في الحلقات.
    \end{columns}
\end{frame}

\begin{frame}{الكتابة على الشاشة}
    \begin{center}
    \begin{minipage}{0.62\textwidth}
        \ccode{d1/write.c}
        \out{d1/write.out}
    \end{minipage}
    \end{center}
    \vspace{0.3cm}
    \code{\%d} تُستبدل بقيمة المتغير، و\code{\textbackslash n} تعني سطرًا جديدًا.
\end{frame}

\begin{frame}{القراءة من لوحة المفاتيح}
    \begin{center}
    \begin{minipage}{0.62\textwidth}
        \ccode{d1/read.c}
    \end{minipage}
    \end{center}
    \vspace{0.3cm}
    \begin{obox}{\faExclamationTriangle\quad انتبه}
        في \code{scanf} نضع \code{\&} قبل اسم المتغير.
    \end{obox}
\end{frame}

\begin{frame}{العمليات الحسابية}
    \begin{center}
    \begin{tabular}{|c|c|c|}
        \hline
        \ras{العملية} & \ras{\textenglish{C}} & \ras{مثال} \\ \hline
        جمع، طرح، ضرب & \code{+ - *} & \code{7 * 2} $\to$ \textenglish{14} \\ \hline
        القسمة        & \code{/}     & \code{7.0 / 2} $\to$ \textenglish{3.5} \\ \hline
        القسمة الصحيحة & \code{/}     & \code{7 / 2} $\to$ \textenglish{3} \\ \hline
        باقي القسمة   & \code{\%}    & \code{7 \% 2} $\to$ \textenglish{1} \\ \hline
    \end{tabular}
    \end{center}
\end{frame}

\begin{frame}{القسمة بين عددين صحيحين}
    \begin{columns}[T]
        \column{0.55\textwidth}
        \ccode{d1/div.c}
        \column{0.4\textwidth}
        \out{d1/div.out}
    \end{columns}
    \vspace{0.3cm}
    \code{int / int} تحذف الجزء العشري. نكتب \code{(float)} للحصول على نتيجة حقيقية.
\end{frame}

\begin{frame}{مثال كامل: مساحة دائرة}
    \begin{center}
    \begin{minipage}{0.62\textwidth}
        \ccode{d1/circle.c}
    \end{minipage}
    \end{center}
\end{frame}

\begin{frame}{تنفيذ البرنامج}
    \vspace{0.5cm}
    \begin{center}
        \begin{minipage}{0.5\textwidth}
            \out{d1/circle.out}
        \end{minipage}
    \end{center}
    \vspace{0.3cm}
    يكتب المستخدم \textenglish{2}، فيعرض البرنامج المساحة برقمين بعد الفاصلة (\code{\%.2f}).
\end{frame}

\begin{frame}{جدول التتبع: تبديل قيمتين}
    \begin{columns}[T]
        \column{0.38\textwidth}
        \ccode{d1/swap.c}
        \column{0.55\textwidth}
        \begin{center}
        \begin{tabular}{|c|c|c|c|}
            \hline
            \ras{السطر} & \ras{\code{a}} & \ras{\code{b}} & \ras{\code{t}} \\ \hline
            1 & 5 & 8 & -- \\ \hline
            2 & 5 & 8 & \cell{tealDeep}{\textbf{5}} \\ \hline
            3 & \cell{tealDeep}{\textbf{8}} & 8 & 5 \\ \hline
            4 & 8 & \cell{tealDeep}{\textbf{5}} & 5 \\ \hline
        \end{tabular}
        \end{center}
    \end{columns}
    \vspace{0.2cm}
    نحتاج إلى متغير مساعد \code{t}، لأن التعيين يمحو القيمة القديمة.
\end{frame}

\begin{frame}{أمثلة تطبيقية}
    \begin{tbox}{\faLaptopCode\quad عشرة أمثلة من مجالات مختلفة}
        نطبق ما تعلمناه في مجالات من الحياة اليومية:
        الصحة، الرياضة، الطقس، النقل، التجارة الإلكترونية، الهواتف الذكية، الشبكات الاجتماعية، الطاقة، البيئة، البث الرقمي.
    \end{tbox}
\end{frame}

\begin{frame}[t]{المثال الأول: الصحة}
    \begin{columns}[T]
        \column{0.64\textwidth}
        \ccode[basicstyle=\codefont\codesmall]{d1/ex01.c}
        \column{0.33\textwidth}
        {\small\textbf{مؤشر كتلة الجسم}: نحسب مؤشر كتلة الجسم: الوزن مقسوم على مربع الطول.\par}
        \vspace{0.3cm}
        \out[basicstyle=\codefont\scriptsize\color{brandDark}]{d1/ex01.out}
    \end{columns}
\end{frame}

\begin{frame}[t]{المثال الثاني: الرياضة}
    \begin{columns}[T]
        \column{0.64\textwidth}
        \ccode[basicstyle=\codefont\codesmall]{d1/ex02.c}
        \column{0.33\textwidth}
        {\small\textbf{سرعة عدّاء}: عدّاء يقطع 10 كم في 52 دقيقة: نحسب سرعته ووتيرته.\par}
        \vspace{0.3cm}
        \out[basicstyle=\codefont\scriptsize\color{brandDark}]{d1/ex02.out}
    \end{columns}
\end{frame}

\begin{frame}[t]{المثال الثالث: الطقس}
    \begin{columns}[T]
        \column{0.64\textwidth}
        \ccode[basicstyle=\codefont\codesmall]{d1/ex03.c}
        \column{0.33\textwidth}
        {\small\textbf{تحويل درجة الحرارة}: نحوّل درجة الحرارة من سلسيوس إلى فهرنهايت.\par}
        \vspace{0.3cm}
        \out[basicstyle=\codefont\scriptsize\color{brandDark}]{d1/ex03.out}
    \end{columns}
\end{frame}

\begin{frame}[t]{المثال الرابع: النقل}
    \begin{columns}[T]
        \column{0.64\textwidth}
        \ccode[basicstyle=\codefont\codesmall]{d1/ex04.c}
        \column{0.33\textwidth}
        {\small\textbf{استهلاك الوقود}: نحسب استهلاك سيارة في 100 كم وتكلفة الرحلة.\par}
        \vspace{0.3cm}
        \out[basicstyle=\codefont\scriptsize\color{brandDark}]{d1/ex04.out}
    \end{columns}
\end{frame}

\begin{frame}[t]{المثال الخامس: التجارة الإلكترونية}
    \begin{columns}[T]
        \column{0.64\textwidth}
        \ccode[basicstyle=\codefont\codesmall]{d1/ex05.c}
        \column{0.33\textwidth}
        {\small\textbf{السعر بعد التخفيض}: نطبق تخفيضًا بنسبة مئوية ثم نضيف ثمن التوصيل.\par}
        \vspace{0.3cm}
        \out[basicstyle=\codefont\scriptsize\color{brandDark}]{d1/ex05.out}
    \end{columns}
\end{frame}

\begin{frame}[t]{المثال السادس: الهواتف الذكية}
    \begin{columns}[T]
        \column{0.64\textwidth}
        \ccode[basicstyle=\codefont\codesmall]{d1/ex06.c}
        \column{0.33\textwidth}
        {\small\textbf{مدة البطارية}: نحسب عدد ساعات عمل الهاتف من سعة البطارية والاستهلاك.\par}
        \vspace{0.3cm}
        \out[basicstyle=\codefont\scriptsize\color{brandDark}]{d1/ex06.out}
    \end{columns}
\end{frame}

\begin{frame}[t]{المثال السابع: الشبكات الاجتماعية}
    \begin{columns}[T]
        \column{0.64\textwidth}
        \ccode[basicstyle=\codefont\codesmall]{d1/ex07.c}
        \column{0.33\textwidth}
        {\small\textbf{معدل التفاعل}: نسبة التفاعل = (الإعجابات + التعليقات) مقسومة على المتابعين.\par}
        \vspace{0.3cm}
        \out[basicstyle=\codefont\scriptsize\color{brandDark}]{d1/ex07.out}
    \end{columns}
\end{frame}

\begin{frame}[t]{المثال الثامن: الطاقة}
    \begin{columns}[T]
        \column{0.64\textwidth}
        \ccode[basicstyle=\codefont\codesmall]{d1/ex08.c}
        \column{0.33\textwidth}
        {\small\textbf{فاتورة الكهرباء}: مدفأة استطاعتها 1500 واط تعمل 6.5 ساعات: نحسب الطاقة والتكلفة.\par}
        \vspace{0.3cm}
        \out[basicstyle=\codefont\scriptsize\color{brandDark}]{d1/ex08.out}
    \end{columns}
\end{frame}

\begin{frame}[t]{المثال التاسع: البيئة}
    \begin{columns}[T]
        \column{0.64\textwidth}
        \ccode[basicstyle=\codefont\codesmall]{d1/ex09.c}
        \column{0.33\textwidth}
        {\small\textbf{البصمة الكربونية}: نحسب انبعاثات سيارة في شهر، وعدد الأشجار اللازمة لامتصاصها في سنة.\par}
        \vspace{0.3cm}
        \out[basicstyle=\codefont\scriptsize\color{brandDark}]{d1/ex09.out}
    \end{columns}
\end{frame}

\begin{frame}[t]{المثال العاشر: البث الرقمي}
    \begin{columns}[T]
        \column{0.64\textwidth}
        \ccode[basicstyle=\codefont\codesmall]{d1/ex10.c}
        \column{0.33\textwidth}
        {\small\textbf{مدة فيديو}: نحوّل مدة فيديو بالثواني إلى ساعات ودقائق وثوانٍ، باستعمال \code{/} و\code{\%}.\par}
        \vspace{0.3cm}
        \out[basicstyle=\codefont\scriptsize\color{brandDark}]{d1/ex10.out}
    \end{columns}
\end{frame}


\begin{frame}{أخطاء شائعة}
    \begin{obox}{\faExclamationTriangle\quad تجنّب هذه الأخطاء}
        \begin{itemize}
            \item نسيان الفاصلة المنقوطة \code{;} أو الرمز \code{\&} في \code{scanf}.
            \item صيغة لا توافق النوع: قراءة \code{float} بالصيغة \code{\%d}.
            \item استعمال متغير قبل إعطائه قيمة.
        \end{itemize}
    \end{obox}
\end{frame}

\begin{frame}{خلاصة الدرس}
    \begin{tbox}{\faCheckCircle\quad خلاصة الدرس}
        \begin{itemize}
            \item الخوارزمية تعليمات منتهية ومرتبة ودقيقة.
            \item المتغير قيمته تتغير، والثابت قيمته لا تتغير.
            \item التعيين \code{=}، القراءة \code{scanf}، الكتابة \code{printf}.
        \end{itemize}
    \end{tbox}
\end{frame}

\end{document}
